\chapter*{Sommario}
\addcontentsline{toc}{chapter}{Sommario}
\markboth{Sommario}{Sommario}

Il progetto europeo UPRISE, nato dopo il sisma del 2016 nell'area di Camerino,
prevede l'integrazione di sensori IoT negli arredi scolastici: in caso di terremoto
il nodo installato nell'arredo deve rilevare la presenza di una persona rifugiata
sotto di esso e comunicarlo all'esterno tramite rete LoRa. Il nodo attualmente in
uso (DIPME-DEVICE) affida il rilevamento di presenza a un sensore \emph{PIR}, che
per principio fisico è cieco verso una persona immobile: proprio la condizione
tipica di chi si trova intrappolato.

Questa tesi affronta la parte di \emph{sensing} del problema. Sono stati studiati,
strumentati e confrontati sperimentalmente due radar \emph{mmWave} FMCW a 24~GHz
(HLK-LD2410B e HLK-LD2420) e un sensore PIR HC-SR501, acquisendo i dati con una
piattaforma di misura realizzata su ESP32 che campiona a 5~Hz in
\emph{engineering mode}, cioè con accesso ai valori di energia riflessa per ciascuno
dei nove \emph{gate} di distanza del radar.

Il risultato centrale è quantitativo: su cinque prove indipendenti da cinque minuti
ciascuna, con un soggetto immobile seduto a 2,30~m, il PIR ha mancato la presenza
nel $99{,}78 \pm 0{,}49\,\%$ dei campioni, mentre il radar l'ha rilevata nel
$100\,\%$ dei 7554 campioni acquisiti. Il PIR risulta inaffidabile anche con il
soggetto \emph{in movimento sul posto} ($80{,}5\,\%$ di falsi negativi a 1~m,
$100\,\%$ oltre i 2~m), coerentemente con il fatto che l'elemento piroelettrico
risponde al transito del flusso infrarosso fra le zone della lente di Fresnel e non
al movimento in sé. La misura di distanza del radar si è rivelata estremamente
lineare ($R^2 = 0{,}99965$ su 25 prove fra 1 e 5~m), con un errore di scala del
$+3{,}81\,\%$ correggibile mediante un singolo coefficiente; su 7~ore di stanza
vuota nessuno dei due sensori ha prodotto falsi positivi, con limite superiore al
95\,\% di confidenza pari a $0{,}43$ eventi/h. Dall'analisi in frequenza
dell'energia per-gate è stato inoltre possibile estrarre il ritmo respiratorio di un
soggetto immobile, con concordanza fra canali indipendenti.

Su questa base la tesi propone un \emph{indice di vitalità}: una grandezza compresa
fra 0 e 100 che, oltre al ``c'è qualcuno'', stimi quanto la persona si muove, a
supporto del triage dei soccorritori. Il lavoro comprende infine l'analisi dei
consumi da datasheet e il progetto di un'interfaccia web servita direttamente
dall'ESP32 per la visualizzazione in tempo reale e l'esportazione dei dati.

La raccomandazione conclusiva non è la sostituzione del PIR, ma
un'\emph{architettura ibrida}: il PIR come guardiano a microampere che decide quando
accendere il radar, e il radar mmWave come strumento di misura durante la finestra
di emergenza.

\vspace{1em}
\begin{sloppypar}
\noindent\textbf{Parole chiave:} radar mmWave, FMCW 24~GHz, PIR, rilevamento di
presenza, ESP32, HLK-LD2410B, monitoraggio del respiro, sicurezza sismica,
UPRISE/SAFE.
\end{sloppypar}
